\chapter*{Introduction}
\addcontentsline{toc}{chapter}{Introduction} % Ajoute l'intro à la table des matières

Depuis l'introduction de l'architecture Transformer \cite{vaswani2017attention}, nous avons assisté à l'arrivée de modèles de langage de plus en plus performants. Initialement limités (comme GPT-2 ou BERT), ces modèles démontrent aujourd'hui des capacités linguistiques comparables, voire supérieures, à celles des humains sur certaines tâches spécifiques. Cependant, malgré ces performances, ces réseaux de neurones demeurent souvent des « boîtes noires » dont les mécanismes internes d'émergence du sens restent méconnus.

Dans le cadre de ce projet tuteuré, nous nous intéressons à la comparaison entre le fonctionnement de ces modèles artificiels et les mécanismes neurobiologiques du cerveau humain. Notre problématique centrale consiste à déterminer si l'Intelligence Artificielle traite la sémantique de manière analogue au cerveau, particulièrement en ce qui concerne l'unification des mots pour former une phrase cohérente. Pour ce faire, nous nous appuyons sur le modèle neurolinguistique MUC (\textit{Memory, Unification, Control}) \cite{hagoort2013memory} ainsi que sur une métrique issue de la géométrie de l'information : la dimensionnalité intrinsèque (ID) \cite{facco2017estimating, ansuini2019intrinsic}.

Pour mener à bien cette étude, nous avons choisi d'implémenter un environnement expérimental en Python, s'appuyant sur la librairie PyTorch. Nous utilisons le modèle GPT-2 pour son architecture causale (\textit{decoder-only}), proche de la linéarité de la lecture humaine. Nous utiliserons l'algorithme TwoNN \cite{facco2017estimating} pour estimer la complexité des représentations vectorielles (\textit{embeddings}) générées par le modèle. Afin d'isoler le traitement sémantique, nous avons élaboré un protocole comparant le traitement de phrases naturelles à celui de phrases de type « Jabberwocky » (syntaxiquement correctes mais dénuées de sens).

Ce rapport constitue la première partie de notre projet, focalisée sur l'étude théorique et la planification. Dans le premier chapitre, nous dressons un état de l'art reliant l'architecture des Transformers aux théories neurolinguistiques via le concept de dimensionnalité intrinsèque \cite{recanatesi2021predictive}. Dans le chapitre suivant, nous détaillons notre méthodologie et justifions nos choix d'implémentation. Enfin, nous présentons dans le troisième chapitre la feuille de route (\textit{roadmap}) détaillée de la phase d'expérimentation qui se déroulera au second semestre.